(1) 纯态关于\(\phi\)的量子Fisher信息为\(F_Q=4[\langle\partial_\phi\psi|\partial_\phi\psi\rangle-|\langle\psi|\partial_\phi\psi\rangle|^2]\)。对单光子态，\(|\partial_\phi\psi_1\rangle=ie^{i\phi}|0,1\rangle/\sqrt2\)，前项为\(1/2\)、内积模方为\(1/4\)，所以\(F_Q^{(1)}=4(1/2-1/4)=1\)。
(2) 对\(N\)次独立同分布单光子实验，量子Fisher信息相加\(F_Q^{\rm prod}=N\)。量子Cramér–Rao界给无偏估计方差\(\operatorname{Var}(\hat\phi)\ge1/(\nu N)\)，其中\(\nu\)为独立重复实验次数；一次配置的标度\(\Delta\phi\gtrsim1/\sqrt N\)就是散粒噪声标准量子极限。
(3) NOON态导数\(|\partial_\phi\psi_N\rangle=iNe^{iN\phi}|0,N\rangle/\sqrt2\)，故\(\langle\partial\psi|\partial\psi\rangle=N^2/2\)、\(|\langle\psi|\partial\psi\rangle|^2=N^2/4\)，得到\(F_Q^{\rm NOON}=N^2\)，理想一次配置界为\(\Delta\phi\gtrsim1/N\)。增强来源于累积\(N\phi\)的多光子相位而非违反量子测量统计。
(4) 若仅保留所有\(N\)个光子均存活的后选择事件，其成功率\(p_{\rm surv}=\eta^N\)，成功时仍有\(F_Q=N^2\)，所以按每次原始制备计算的有效信息为\(F_{\rm eff}=p_{\rm surv}N^2=\eta^N N^2\)。对固定\(\eta<1\)，随\(N\)增长此因子最终指数衰减，表明理想Heisenberg标度对损耗十分敏感；实际损耗后的非后选择QFI需按混态模型另算。
